\chapter{结论与展望}

\section{结论}

作为一项前沿应用技术，人工智能大模型已广泛融入众多领域，为人们的生产生活带来极大便利。随着人工智能大模型的迅猛发展，各技术公司面临的相关开发测试任务日益繁重。为顺应AI发展的实际需求，本论文聚焦于AI多任务调度问题的特性展开深入研究，着力解决现有编排模式与调度系统无法满足AI大模型任务发展需求的难题。在此基础上，本系统精心设计了三个工程层次，并对各层次内部进行了细致规划。具体研究工作与贡献如下：

\begin{enumerate}
\item 攻克了传统调度系统对AI相关特性适配不佳的难题。API 功能接入层与任务执行层紧密围绕人工智能大模型的独特需求展开深度适配。在参数类型方面，充分考虑到 AI 大模型处理数据的多样性和复杂性，对参数进行了精细化设计，确保能够准确传递各种类型的数据，满足不同模型的输入要求。在数据传输方式上，针对 AI 大模型数据量大、实时性要求高的特点，选用了合适的传输协议和策略，减少数据传输延迟和丢包率，实现了数据的快速、稳定传输。通过这些维度的深度适配，使系统在各个环节高度贴合人工智能大模型任务的特性，有力地化解了现有任务调度系统与人工智能大模型任务适配度不理想的困境，为 AI 任务的高效执行奠定了坚实基础。
\item 成功解决了现阶段无法准确判断 AI 任务是否具备执行条件的关键问题。任务调度层创新性地设计了任务可用性检测机制，并将其与先进的调度算法深度融合、协同运作。该机制基于对 AI 任务的详细分析，综合考虑任务的优先级、资源需求、前置条件以及当前系统的资源状态等多方面因素，构建了一套判断体系。借助该机制与算法的协同作用，本系统能够精准洞察 AI 任务是否具备执行条件，避免了因盲目执行任务而导致的资源浪费和效率低下问题。这一举措不仅有效节省了大量的人力成本，还极大地提高了系统的整体运行效率和稳定性，确保了 AI 任务在复杂多变的环境中能够有条不紊地推进。
\item 设计了专门针对 AI 特有资源的资源互斥检测机制，解决了在 AI 多任务场景下难以识别资源互斥的问题。任务调度层不仅设计了资源可用性检测机制，还结合 AI 多任务场景下资源使用的特点，开发了一套针对性的检测算法。该算法能够实时监测各个任务对资源的占用情况，准确识别出 AI 多任务场景下潜在的资源互斥状况。收到任务时，若检测到资源互斥，系统能够迅速做出响应，通过合理调整任务执行顺序或重新分配资源等策略，避免因资源冲突而导致的任务失败或系统性能下降。这一机制的引入显著提高了调度效率，保障了 AI 多任务在共享资源环境下的高效、稳定运行。
\item 顺利完成了对本系统的全面整合测试。测试阶段，采用了多种测试方法和场景模拟，涵盖功能测试、性能测试、稳定性测试等多个维度。经过大量的实际测试和优化，本系统在实际应用中展现出卓越的性能。
\end{enumerate}

\section{展望}

本系统的设计与实现能够满足当前腾讯公司AI任务调度的需求，但AI大模型作为一个新兴领域，目前仍处于发展初期。在可预见的未来，AI大模型将持续快速演进。因此，即便本系统当前能够适配现有需求，仍需紧跟技术发展趋势，不断进行调整与优化，才能在AI技术的迭代浪潮中保持领先优势。本系统需要继续拓展与挖掘的地方如下：

\begin{enumerate}
    \item 适配更多样的临界资源。AI发展的早期，其依赖的临界资源主要是GPU服务器。而发展至今日，AI需求的临界资源早已不局限于硬件资源。以最近热议的MCP为例，有的MCP依赖的资源也是互斥使用的，本系统可以以此为创新点，挖掘与MCP服务对接的可能性。
    \item 完善数据报表功能。当前数据报表的展示格式主要由开发者与系统维护人员沟通实现。日后可以增加AI大模型自动分析任务数据，自动生成报表的功能。在基础数据之上，利用AI智能分析数据报表，既可以节省沟通成本，又可以从不同的视角分析数据。
\end{enumerate}

\makeatletter
\@openrightfalse
\makeatother